\section{Impulse theorem method}

The impulse theorem converts a $\delta(t-\tau)$ forcing into a jump of the initial data at $t=\tau+0$ for the corresponding homogeneous equation. One then solves that Cauchy (or Fourier) problem by known methods.

\subsection{Infinite string with zero initial data}

Consider
\begin{equation}
u_{tt}-a^2 u_{xx}=f(x,t),
\qquad
u\big|_{t=0}=u_t\big|_{t=0}=0.
\end{equation}
The Green's function satisfies the homogeneous wave equation for $t>\tau$, together with
\begin{equation}
G\big|_{t=\tau+0}=0,
\qquad
G_t\big|_{t=\tau+0}=\delta(x-\xi).
\end{equation}
d'Alembert's formula with the time shift $t\mapsto t-\tau$ gives
\begin{equation}
G(x,t;\xi,\tau)
=
\begin{cases}
0, & \xi\notin\bigl(x-a(t-\tau),\,x+a(t-\tau)\bigr),\\[0.3em]
\dfrac{1}{2a}, & \xi\in\bigl(x-a(t-\tau),\,x+a(t-\tau)\bigr).
\end{cases}
\end{equation}
Hence
\begin{equation}
u(x,t)
=\int_0^{t}\int_{x-a(t-\tau)}^{x+a(t-\tau)}
\frac{1}{2a}\,f(\xi,\tau)\,\mathrm{d}\xi\,\mathrm{d}\tau.
\end{equation}

\subsection{Finite interval with Neumann ends}

As a second illustration take a forced string on $[0,\ell]$ with free ends and vanishing initial data,
\begin{equation}
\begin{gathered}
u_{tt}-a^2 u_{xx}=A\cos\frac{\pi x}{\ell}\sin\omega t,\\
u_x\big|_{x=0}=u_x\big|_{x=\ell}=0,
\qquad
u\big|_{t=0}=u_t\big|_{t=0}=0.
\end{gathered}
\end{equation}
The impulse problem is the homogeneous wave equation with $G|_{t=\tau+0}=0$, $G_t|_{t=\tau+0}=\delta(x-\xi)$, and Neumann boundary conditions. Separation of variables yields
\begin{equation}
\begin{aligned}
G(x,t;\xi,\tau)
&=\frac{t-\tau}{\ell}
+\frac{2}{\pi a}\sum_{n=1}^{\infty}\frac{1}{n}
\sin\frac{n\pi a(t-\tau)}{\ell}
\cos\frac{n\pi\xi}{\ell}\cos\frac{n\pi x}{\ell}.
\end{aligned}
\end{equation}
Inserting into $u=\int_0^t\int_0^\ell f\,G\,\mathrm{d}\xi\,\mathrm{d}\tau$ and using orthogonality retains only the $n=1$ mode:
\begin{equation}
u(x,t)
=\frac{A\ell}{\pi a}\frac{1}{\omega^2-\pi^2 a^2/\ell^2}
\Bigl(\omega\sin\frac{\pi a t}{\ell}-\frac{\pi a}{\ell}\sin\omega t\Bigr)
\cos\frac{\pi x}{\ell}.
\end{equation}
When $\omega=\pi a/\ell$ the denominator vanishes and a resonant solution with secular growth must be constructed separately.

\subsection{Infinite line heat equation}

For
\begin{equation}
u_t-a^2 u_{xx}=f,
\qquad
u\big|_{t=0}=0,
\end{equation}
the impulse problem is $G_t-a^2 G_{xx}=0$ with $G|_{t=\tau+0}=\delta(x-\xi)$. The heat kernel is
\begin{equation}
\label{eq:heat-kernel}
G(x,t;\xi,\tau)
=\frac{1}{2a\sqrt{\pi(t-\tau)}}
\exp\Bigl(-\frac{(x-\xi)^2}{4a^2(t-\tau)}\Bigr),
\end{equation}
and therefore
\begin{equation}
u(x,t)
=\int_0^{t}\int_{-\infty}^{\infty}
f(\xi,\tau)\,G(x,t;\xi,\tau)\,\mathrm{d}\xi\,\mathrm{d}\tau.
\end{equation}

\begin{note}[Infinite propagation speed]
As a function of $x-\xi$, the kernel \eqref{eq:heat-kernel} is a Gaussian that widens and flattens as $t-\tau$ increases, with its peak at the source. For any $t>\tau$ one has $G\neq0$ everywhere: a signal propagates with infinite speed. This is an artifact of Fourier's law, but the description remains accurate when $t-\tau$ is not extremely small.
\end{note}

\subsection{Semi-infinite heat problems by odd/even extension}

\paragraph{Dirichlet end $u|_{x=0}=0$.}
Extend $f$ oddly to $x<0$. Then \eqref{eq:heat-kernel} produces
\begin{equation}
\begin{aligned}
G(x,t;\xi,\tau)
&=\frac{1}{2a\sqrt{\pi(t-\tau)}}
\Biggl[
\exp\Bigl(-\frac{(x-\xi)^2}{4a^2(t-\tau)}\Bigr)\\
&\qquad\qquad
-\exp\Bigl(-\frac{(x+\xi)^2}{4a^2(t-\tau)}\Bigr)
\Biggr],
\end{aligned}
\end{equation}
and $u=\int_0^t\int_0^{\infty}f(\xi,\tau)G\,\mathrm{d}\xi\,\mathrm{d}\tau$.

\paragraph{Neumann end $u_x|_{x=0}=0$.}
The even extension yields
\begin{equation}
\begin{aligned}
G(x,t;\xi,\tau)
&=\frac{1}{2a\sqrt{\pi(t-\tau)}}
\Biggl[
\exp\Bigl(-\frac{(x-\xi)^2}{4a^2(t-\tau)}\Bigr)\\
&\qquad\qquad
+\exp\Bigl(-\frac{(x+\xi)^2}{4a^2(t-\tau)}\Bigr)
\Biggr].
\end{aligned}
\end{equation}
